% =============================================================================
%  «ТЕОРИЯ СИСТЕМ» — ствол серии «Архитектура Рассуждения»
%  Блок 3. КАТЕГОРИЯ: системы складываются.  Разделы 3.1–3.9 (черновик прозы).
%  Стиль: «выводим и показываем»; репо-якоря — в приложении, не в сносках.
%  Самокомпилируемо: pdflatex blok-3-kategoriya.tex (см. память latex-build-windows).
% =============================================================================
\documentclass[12pt,a4paper]{article}
\usepackage{cmap}
\usepackage[T2A]{fontenc}
\usepackage[utf8]{inputenc}
\usepackage[russian]{babel}
\usepackage{paratype}
\usepackage{amsmath,amssymb,amsthm}
\usepackage{listings}
\usepackage[expansion=false]{microtype}
\usepackage[margin=28mm]{geometry}

\newcommand{\rezultat}{\medskip\noindent\emph{Что отсюда следует.}\ }

\title{\textbf{Теория Систем}\\[2mm]\large Блок 3. Категория: системы складываются}
\author{}
\date{}

\begin{document}
\maketitle
\thispagestyle{empty}

\noindent\small\textit{Черновик: разделы 3.1--3.9. Продолжает Блок~2 (объект). Категория здесь ---
не общий язык математики (это Том «Математика»), а категория \emph{именно} систем-триад. Машинные
якоря --- в Приложении в конце блока.}\normalsize

\bigskip

\noindent Объект разобран изнутри --- покажем теперь, как объекты \emph{складываются}. Системы
образуют \emph{категорию}: между ними есть отображения, которые можно компоновать, и в этой
категории работают привычные конструкции --- произведение, фактор, изоморфизм, конечные пределы.
Её собственная черта: морфизм \emph{сохраняет роли}, а складывать можно там, где \emph{правила это
пропускают}.

% =============================================================================
\section*{3.1\quad Морфизм; правила гейтят композицию}
% =============================================================================

Морфизм систем --- отображение носителя одной в носитель другой, \emph{сохраняющее роли}: что было
соотнесено, остаётся соотнесённым. Правила при этом не переносятся --- они \emph{рамка}, которая
ограничивает, что вообще составимо. Есть тождественный морфизм и композиция; законы категории
(ассоциативность, единицы) выполняются просто потому, что это композиция функций.

Гейт правил виден сразу. Системы-\emph{эквивалентности} составляются в произведение, а
система-\emph{линейный порядок} --- нет: две несравнимые точки не дают согласованной соотнесённости.
Правила --- твёрдый ярус; это продолжение ранговой асимметрии (§2.4) уже на композиции.

\rezultat
Системы --- \textbf{категория}: морфизм хранит роли, а правила \emph{гейтят} составимость. Складывание
не свободно --- оно проходит там, где правила пропускают.

% =============================================================================
\section*{3.2\quad Произведение и копроизведение}
% =============================================================================

Произведение двух систем --- \emph{настоящее} категорное: есть проекции на сомножители, и всякий
согласованный морфизм в обе системы пропускается через \emph{единственного} посредника. Это работает
не только для эквивалентностей, а для любой конституции, замкнутой относительно произведения.

Двойственно --- копроизведение, сумма: инъекции из слагаемых и единственный посредник из обеих.

\rezultat
Системы \textbf{умножаются и складываются} с универсальным свойством: посредник всегда \emph{единствен}.
Это первый признак того, что перед нами полноценная категория, а не просто набор отображений.

% =============================================================================
\section*{3.3\quad Под-система и фактор}
% =============================================================================

Двойственная пара операций. \emph{Под-система} --- вложение на конституции, устойчивой к ограничению:
часть остаётся системой. \emph{Фактор} --- огрубление ролей до конгруэнции (эпиморфизм): мы
\emph{склеиваем} то, что система не различает, и получаем систему-частное.

\rezultat
\textbf{Вырезать часть} и \textbf{склеить целое} --- две двойственные операции. Вырезание работает на
носителе (Elements), склейка --- на ролях (Roles); эту асимметрию подхватит §3.6.

% =============================================================================
\section*{3.4\quad Изоморфизм и первая теорема}
% =============================================================================

Изоморфизм систем --- это \emph{моно $+$ эпи}: вложение, которое и однозначно, и накрывает.

Первая теорема об изоморфизме держится в \emph{локализованной} форме. Ядро морфизма конгруэнтно ---
то есть честно задаёт фактор --- \emph{ровно тогда}, когда морфизм «схлопывает роли»; оставшаяся часть
условия --- инъективность посредника. Препятствие здесь \emph{названо явно}, а не спрятано в общую
формулировку.

\rezultat
Изоморфизм $=$ моно $+$ эпи; первая теорема об изоморфизме работает при honest-условии на роли.
\textbf{Локализованная форма --- это точность, а не слабость:} мы показываем, \emph{где} именно проходит
граница её применимости.

% =============================================================================
\section*{3.5\quad Терминал, инициал; нет нулевого объекта}
% =============================================================================

Есть \emph{терминальный} объект --- одноэлементная система, в которую из любой ведёт единственный
морфизм; и \emph{инициальный} --- пустая, из которой в любую ведёт единственный. Но это \emph{разные}
объекты: \emph{нулевого} --- одновременно терминального и инициального --- нет.

\rezultat
Терминал и инициал есть, нулевого нет. Это \textbf{черта set-подобия}, а не дефект: систем-категория
ближе к множествам, чем к группам или модулям, --- и эта близость сама есть содержательный факт.

% =============================================================================
\section*{3.6\quad Конечная биполнота}
% =============================================================================

Категория замкнута относительно \emph{всех конечных} пределов и копределов: эквалайзеры и
пулбэки (пределы), коэквалайзеры и пушауты (копределы).

И здесь --- тонкая двойственная асимметрия, первый отчётливый намёк на границу. Эквалайзер
\emph{вырезает} носитель --- конкретное под-множество, на ярусе Elements; коэквалайзер \emph{сливает}
в роли --- на ярусе Roles. Вырезание всегда конкретно; а слияние носителя --- ровно там, где в
Блоке~7 обнаружится role-limit.

\rezultat
Категория \textbf{конечно биполна}. Но в самой двойственности эквалайзера и коэквалайзера ---
вырезать против слить --- уже слышен звон единственной границы теории (§4.5, Блок~7).

% =============================================================================
\section*{3.7\quad Регулярность}
% =============================================================================

Всякий морфизм раскладывается \emph{канонически}: сперва эпиморфизм \emph{на образ}, затем
мономорфизм-вложение образа. Образ --- это под-система на тех элементах, которых морфизм достигает.

\rezultat
Морфизм $=$ накрытие образа $\circ$ вложение образа. Категория \textbf{регулярна}: у каждого
отображения есть честный «образ-как-под-система».

% =============================================================================
\section*{3.8\quad Второе измерение: Roles как 2-клетки}
% =============================================================================

Между морфизмами есть ещё один ярус --- \emph{2-клетки}, и их роль играют сами \emph{Roles}:
«одинаково с точки зрения системы» --- это категорный P3. Внутренний hom делает категорию
\emph{обогащённой над собой}: композиция и единица оказываются \emph{настоящими} морфизмами, а не
внешней приправой.

\rezultat
Систем-категория \textbf{двумерна и обогащена над собой}: Roles суть 2-клетки. Честно: 2-клетки
\emph{Prop-тонкие} --- они различают «так или не так», но не несут высших данных.

% =============================================================================
\section*{3.9\quad Категория явно; автоморфизмы; биекции}
% =============================================================================

Всё сказанное --- не метафора: E/R/R \emph{явно} удовлетворяет аксиомам категории (тождество,
композиция, ассоциативность доказаны прямо).

У объектов есть \emph{автоморфизмы} --- обратимые морфизмы; \emph{дискретная} симметрия (перестановки)
\emph{выведена}. Здесь --- честный зазор: \emph{непрерывная} группа $SU(N)$ \emph{постулируется}, ибо
требует топологии, которой в самом различении нет; а число $N^2-1$ --- размерность алгебры Ли, не
счёт дискретных генераторов. Группы автоморфизмов смотрят в сторону калибровки --- это нить Блока~8.

И четыре \emph{биекции} показывают одно: одна структура отвечает на четыре разных вопроса
(процесс $\leftrightarrow$ система, сжатие $\leftrightarrow$ физика,
декогеренция $\leftrightarrow$ затухание, наблюдатель $\leftrightarrow$ компрессор).

\rezultat
Категория --- \textbf{доказанная структура}, а не образ речи; её автоморфизмы --- мост к калибровке,
проведённый \emph{с честным зазором} ($SU(N)$ постулируется, не выводится).

\medskip
\noindent\emph{Переход.} Статичный объект разобран как категория --- биполная, регулярная, двумерная,
обогащённая над собой. Следующий блок даёт ему \emph{время}: оператор, эволюция, две стрелы.

% =============================================================================
\appendix
\section*{Приложение к Блоку 3 — машинные якоря}
% =============================================================================
\small\raggedright\sloppy
\noindent Всё перечисленное машинно проверено; единственная аксиома ветви --- \texttt{classic} (L3).
Кросс-ссылка: общая теория категорий как язык математики --- в Томе «Математика» (Часть~XIV); здесь ---
категория \emph{именно} \texttt{FunctionalSystem}.

\begin{itemize}\itemsep2pt
\item \textbf{§3.1.} \texttt{ERRMorphism} (сохраняет Roles), \texttt{err\_id\_left/right},
\texttt{err\_comp\_assoc}; гейт --- \texttt{equiv\_product\_closed} (да) /
\texttt{connex\_not\_product\_closed} (нет).\\
\emph{Файл:} \texttt{ERRComposition.v}.

\item \textbf{§3.2.} \texttt{fs\_product}, общее \texttt{ERRProductGeneral},
\texttt{product\_mediator\_unique}; копроизведение \texttt{sum\_rel}.\\
\emph{Файлы:} \texttt{ERRProductUniversal.v}, \texttt{ERRProductGeneral.v}, \texttt{ERRCoproduct.v}.

\item \textbf{§3.3.} под-система \texttt{fs\_subsystem}/\texttt{fs\_incl}
(\texttt{ERRTierIIResidue.v}); фактор \texttt{fs\_quot} (\texttt{ERRQuotient.v}).

\item \textbf{§3.4.} \texttt{iso} $=$ mono$+$epi (\texttt{ERRIso.v}); \texttt{kernel},
\texttt{roles\_collapsing} (\texttt{ERRFirstIso.v}).

\item \textbf{§3.5.} терминал/инициал, нет нулевого --- \texttt{ERRTerminalInitial.v}.

\item \textbf{§3.6.} \texttt{ERREqualizer.v}, \texttt{ERRCoequalizer.v} (конечная биполнота;
двойственность вырезать\,/\,слить).

\item \textbf{§3.7.} регулярность (эпи-на-образ $\circ$ моно) --- \texttt{ERRImageFactorization.v}.

\item \textbf{§3.8.} \texttt{ERR2Category.v}, \texttt{ERR2CategoryHom.v},
\texttt{ERREnrichedComposition.v} (Roles $=$ 2-клетки; обогащение над собой).

\item \textbf{§3.9.} \texttt{ERRCategory.v} (\texttt{err\_category\_synthesis});
\texttt{ERRAutomorphism.v} (\texttt{err\_automorphism\_synthesis}; \ensuremath{\boldsymbol{\triangle}}\,$SU(N)$ постулируется);
\texttt{ERRBijections.v} (\texttt{four\_bijections}).
\end{itemize}
\normalsize

% --- Дальше: Блок 4 (ВРЕМЯ) по тезисному плану в Книги/Теория Систем/. ---

\end{document}
